\subsection{Fases da manobra do manipulador robótico}
\label{sec:fases_manobra_MR}

A manobra do \gls{mr} é decomposta em quatro fases operacionais, definidas em função da posição relativa na região de captura e do estado de atitude da base:

\begin{enumerate}[label=(\roman*)]
\item \textit{Posicionamento inicial}: o manipulador parte de uma configuração travada e é conduzido a uma postura de trabalho, afastando a ferramenta da base física e orientando os elos em direção à porta de acoplamento.

\item \textit{Alinhamento do efetuador}: a seguir, o manipulador é comandado para alinhar o efetuador com o eixo da porta. Nesta fase, pequenos ajustes de configuração das juntas são usados para reduzir erros laterais e angulares, mantendo a base dentro dos critérios de atitude estabelecidos.

\item \textit{Aproximação final e inserção}: uma vez obtido o alinhamento desejado, a ferramenta avança ao longo do eixo da porta de acoplamento com velocidades reduzidas para realizar a inserção de forma controlada evitando impactos bruscos.

\item \textit{Amortecimento de erros residuais}: após a inserção, são realizados pequenos movimentos de correção para eliminar folgas, reduzir esforços residuais e estabilizar a configuração final do manipulador e da base.
\end{enumerate}